\begin{figure}[htbp]
  \centering
  \includegraphics[width=160mm]{q2_posterior.pdf}
  \caption{首次正常示向度观测前后的位置信息（\(a=0,\beta=0\)）。右图按源距和相对示向角展开狭窄角域，颜色仍表示单位面积上的位置概率密度。}
  \label{fig:q2-posterior}
\end{figure}

\begin{figure}[htbp]
  \centering
  \includegraphics[width=160mm]{q2_feedback.pdf}
  \caption{第二次反馈的预测分布与剩余区域（\(a=0,\beta=0,w=0.5\)，\(q=(813.82,521.02)\) 米）。本例强信号与无信号概率均为0，阴影面积即完成概率；两个下图采用相同长度尺度，对比剩余区域能否被20米圆覆盖。}
  \label{fig:q2-feedback}
\end{figure}

\begin{figure}[htbp]
  \centering
  \includegraphics[width=160mm]{q2_candidates.pdf}
  \caption{11个权重下的候选检测点及评价指标（\(a=0,\beta=0\)）。颜色表示完成概率权重，虚线仅连接相邻权重的数值结果，不表示点间所有位置都已入选。}
  \label{fig:q2-candidates}
\end{figure}
